% ============================================================
\chapter{Bayesian Paradigm}
\label{ch:paradigm}
% ============================================================

In 1763, two years after his death, the Reverend Thomas Bayes saw his essay
published posthumously by the Royal Society of London. The paper proposed a
method for inverting probabilities: given observed data, what can we say
about the cause that produced them? This seemingly innocent question would
spark one of the longest-running debates in all of statistics. On one side,
the \emph{frequentists}, for whom parameters are fixed unknown constants. On
the other, the \emph{Bayesians}, for whom parameters are random variables
endowed with distributions reflecting our state of knowledge. This chapter
presents the Bayesian paradigm: its logic, its strengths, and the reasons
behind its spectacular modern revival.

\begin{intuition}[title=\textbf{Central idea}]
The Bayesian paradigm treats the parameter $\theta$ as a random variable
and encodes all uncertainty through probability distributions. Inference
consists of moving from the prior $\pi(\theta)$ to the posterior
$\pi(\theta \mid x)$ via Bayes' theorem.
\end{intuition}

% ------------------------------------------------------------
\section{Introduction and motivation}
% ------------------------------------------------------------

The frequentist approach regards the parameter $\theta$ as a fixed
unknown constant and evaluates inference procedures by their repeated-sampling
properties. The Bayesian approach, by contrast, models $\theta$ as a
random variable whose distribution reflects our state of knowledge.

This philosophical difference has profound consequences:
\begin{itemize}
  \item Uncertainty is directly quantified by the posterior distribution.
  \item Prior information is formally incorporated.
  \item Any optimal decision can be derived from the posterior.
\end{itemize}

% ------------------------------------------------------------
\section{Axiomatic foundations}
% ------------------------------------------------------------

\subsection{Subjective probability}

\begin{definition}[Subjective probability]
A \textbf{subjective probability} is a function
$\PP : \mathcal{F} \to [0,1]$ satisfying Kolmogorov's axioms,
interpreted as the degree of belief of a rational agent in the
occurrence of an event.
\end{definition}

De Finetti's coherence axioms show that an agent whose bets are
coherent (no \textit{Dutch book}) necessarily behaves as if using
probabilities satisfying Kolmogorov's axioms.

\begin{theorem}[De Finetti's representation theorem]
If $(X_1, X_2, \ldots)$ is an exchangeable sequence of $\{0,1\}$-valued
random variables, then there exists a probability measure $\mu$ on
$[0,1]$ such that:
\[
  \PP(X_1 = x_1, \ldots, X_n = x_n)
  = \int_0^1 \theta^{s_n}(1-\theta)^{n - s_n}\,d\mu(\theta),
\]
where $s_n = \sum_{i=1}^n x_i$.
\end{theorem}

\begin{remark}
This theorem justifies the Bayesian model: exchangeability (a weaker
assumption than independence) implies the existence of a latent
parameter $\theta$ with a prior distribution $\mu$.
\end{remark}

\subsection{Savage's axioms}

\begin{theorem}[Savage's theorem, 1954]
Under Savage's seven axioms (ordering, sure-thing principle, etc.),
there exist:
\begin{enumerate}
  \item a subjective probability measure $\PP$ on the state space,
  \item a utility function $u$ unique up to affine transformation,
\end{enumerate}
such that the agent prefers action $a$ to action $b$ if and only if
$\E_\PP[u(a)] > \E_\PP[u(b)]$.
\end{theorem}

% ------------------------------------------------------------
\section{Bayes' theorem}
% ------------------------------------------------------------

\begin{theorem}[Bayes' theorem]
\label{thm:bayes}
Let $\theta \in \Theta$ be a parameter with prior $\pi(\theta)$ and
$x = (x_1, \ldots, x_n)$ a sample with density $f(x \mid \theta)$.
The posterior distribution is:
\[
  \pi(\theta \mid x)
  = \frac{f(x \mid \theta)\,\pi(\theta)}{m(x)},
  \qquad
  m(x) = \int_\Theta f(x \mid \theta)\,\pi(\theta)\,d\theta.
\]
\end{theorem}

\begin{proof}
By the rule for conditional densities:
\[
  \pi(\theta \mid x)
  = \frac{f(x, \theta)}{f(x)}
  = \frac{f(x \mid \theta)\,\pi(\theta)}
         {\int_\Theta f(x \mid \theta)\,\pi(\theta)\,d\theta}.
\]
\end{proof}

\begin{keyformulas}[title=\textbf{Components of Bayes' theorem}]
\begin{align*}
  \text{Prior} &: \quad \pi(\theta) \\
  \text{Likelihood} &: \quad L(\theta \mid x) = f(x \mid \theta) \\
  \text{Evidence} &: \quad m(x) = \int_\Theta L(\theta \mid x)\,\pi(\theta)\,d\theta \\
  \text{Posterior} &: \quad \pi(\theta \mid x) \propto L(\theta \mid x)\,\pi(\theta)
\end{align*}
\end{keyformulas}

% ------------------------------------------------------------
\section{Choice of the prior distribution}
% ------------------------------------------------------------

\subsection{Types of priors}

\begin{definition}[Informative prior]
An \textbf{informative prior} is a distribution $\pi(\theta)$ that
concentrates its mass on a restricted region of $\Theta$, reflecting
substantial prior knowledge.
\end{definition}

\begin{definition}[Vague (noninformative) prior]
A \textbf{vague} or \textbf{noninformative prior} is a distribution
$\pi(\theta)$ that seeks to let the data dominate the inference.
Examples: uniform prior, Jeffreys prior.
\end{definition}

\begin{definition}[Jeffreys prior]
The \textbf{Jeffreys prior} is defined by:
\[
  \pi_J(\theta) \propto \sqrt{\det I(\theta)},
\]
where $I(\theta)$ is the Fisher information matrix.
\end{definition}

\begin{proposition}[Invariance of Jeffreys prior]
The Jeffreys prior is invariant under reparameterization: if
$\phi = g(\theta)$ is a differentiable bijection, then
$\pi_J(\phi) \propto \sqrt{\det I_\phi(\phi)}$.
\end{proposition}

\begin{proof}
Let $\phi = g(\theta)$. By the change-of-variable rule and the
transformation of Fisher information:
\[
  I_\phi(\phi) = \left(\frac{d\theta}{d\phi}\right)^2 I(\theta),
\]
whence $\sqrt{I_\phi(\phi)} = \abs{\frac{d\theta}{d\phi}} \sqrt{I(\theta)}$,
which is exactly the Jacobian of the change of variable.
\end{proof}

\begin{example}
For $X_1, \ldots, X_n \overset{\text{iid}}{\sim} \mathcal{N}(\mu, \sigma^2)$
with $\sigma^2$ known, the Fisher information for $\mu$ is
$I(\mu) = n/\sigma^2$, a constant. The Jeffreys prior is therefore
$\pi_J(\mu) \propto 1$, i.e., the (improper) uniform prior on $\R$.
\end{example}

\subsection{Reference prior}

\begin{definition}[Bernardo's reference prior]
The \textbf{reference prior} maximizes the expected Kullback--Leibler
information between the prior and the posterior:
\[
  \pi^*(\theta) = \argmax_{\pi} \E_x\left[\KL\big(\pi(\theta \mid x) \| \pi(\theta)\big)\right].
\]
\end{definition}

% ------------------------------------------------------------
\section{Fundamental examples}
% ------------------------------------------------------------

\begin{example}
\textbf{Bernoulli--Beta model.}
Let $X_1, \ldots, X_n \overset{\text{iid}}{\sim} \text{Ber}(\theta)$
with $\theta \sim \text{Be}(a, b)$. Setting $s = \sum x_i$:
\begin{align*}
  \pi(\theta \mid x)
  &\propto \theta^s (1-\theta)^{n-s} \cdot \theta^{a-1}(1-\theta)^{b-1} \\
  &= \theta^{a+s-1}(1-\theta)^{b+n-s-1}.
\end{align*}
Hence $\theta \mid x \sim \text{Be}(a + s, \, b + n - s)$.
\end{example}

\begin{example}
\textbf{Normal--Normal model (known variance).}
Let $X_1, \ldots, X_n \overset{\text{iid}}{\sim} \mathcal{N}(\mu, \sigma^2)$
with $\sigma^2$ known, and $\mu \sim \mathcal{N}(\mu_0, \tau_0^2)$.
The posterior is:
\[
  \mu \mid x \sim \mathcal{N}(\mu_n, \tau_n^2),
\]
where
\[
  \tau_n^2 = \left(\frac{1}{\tau_0^2} + \frac{n}{\sigma^2}\right)^{-1}, \qquad
  \mu_n = \tau_n^2 \left(\frac{\mu_0}{\tau_0^2} + \frac{n\bar{x}}{\sigma^2}\right).
\]
\end{example}

\begin{keyformulas}[title=\textbf{Posterior mean as weighted average}]
\[
  \mu_n = w \cdot \mu_0 + (1-w) \cdot \bar{x},
  \qquad w = \frac{\sigma^2/n}{\tau_0^2 + \sigma^2/n}.
\]
As $n \to \infty$, $w \to 0$ and $\mu_n \to \bar{x}$: the posterior
converges to the maximum likelihood estimator.
\end{keyformulas}

% ------------------------------------------------------------
\section{Asymptotic properties}
% ------------------------------------------------------------

\begin{theorem}[Bernstein--von Mises theorem]
Under regularity conditions, the posterior converges in distribution to
a Gaussian centered on the MLE:
\[
  \pi(\theta \mid x_1, \ldots, x_n)
  \xrightarrow{d}
  \mathcal{N}\!\left(\hat{\theta}_{\text{MLE}},\;
  \frac{1}{n}\,I(\theta_0)^{-1}\right),
\]
where $\theta_0$ is the true parameter value.
\end{theorem}

\begin{remark}
This result shows that, for large samples, the choice of prior has
negligible impact. In other words, the Bayesian and frequentist
approaches converge asymptotically.
\end{remark}

\begin{theorem}[Posterior consistency]
Let $\theta_0$ be the true value. Under regularity conditions
(identifiability, prior support containing $\theta_0$), for every
neighborhood $U$ of $\theta_0$:
\[
  \pi(\theta \in U \mid x_1, \ldots, x_n) \xrightarrow{p} 1
  \quad \text{as } n \to \infty.
\]
\end{theorem}

% ------------------------------------------------------------
\section{Likelihood principle}
% ------------------------------------------------------------

\begin{definition}[Likelihood principle]
Two experiments producing the same likelihood function
$L(\theta \mid x)$ (up to a multiplicative constant) must lead to the
same inference about $\theta$.
\end{definition}

\begin{proposition}
Bayesian inference automatically satisfies the likelihood principle,
since the posterior depends on $x$ only through $L(\theta \mid x)$.
\end{proposition}

\begin{warning}[title=\textbf{Frequentist violation}]
Certain frequentist procedures (such as $p$-values) violate the
likelihood principle because they depend on the sample space and
experimental design, not just the observed data.
\end{warning}

% ------------------------------------------------------------
\section{Python implementation}
% ------------------------------------------------------------

\begin{codeblock}[title=\textbf{Bernoulli--Beta model with PyMC}]
\begin{minted}{python}
import pymc as pm
import arviz as az
import numpy as np

# Simulated data
np.random.seed(42)
n = 50
theta_true = 0.3
data = np.random.binomial(1, theta_true, size=n)

# Bayesian model
with pm.Model() as model_beta_binom:
    # Beta(2, 5) prior
    theta = pm.Beta("theta", alpha=2, beta=5)
    # Likelihood
    y_obs = pm.Bernoulli("y_obs", p=theta, observed=data)
    # MCMC sampling
    trace = pm.sample(2000, tune=1000, random_seed=42)

# Posterior summary
summary = az.summary(trace, var_names=["theta"],
                     hdi_prob=0.95)
print(summary)

# Analytic verification
a_post = 2 + data.sum()
b_post = 5 + n - data.sum()
print(f"Analytic posterior: Be({a_post}, {b_post})")
print(f"Analytic mean: {a_post/(a_post+b_post):.4f}")
\end{minted}
\end{codeblock}

\begin{codeblock}[title=\textbf{Visualizing prior--posterior update}]
\begin{minted}{python}
import matplotlib.pyplot as plt
from scipy import stats

fig, ax = plt.subplots(figsize=(8, 5))
theta_grid = np.linspace(0, 1, 500)

# Prior
ax.plot(theta_grid, stats.beta.pdf(theta_grid, 2, 5),
        label="Prior Be(2, 5)", linestyle="--")
# Analytic posterior
ax.plot(theta_grid, stats.beta.pdf(theta_grid, a_post, b_post),
        label=f"Posterior Be({a_post}, {b_post})")
# Normalized likelihood
s = data.sum()
like = theta_grid**s * (1-theta_grid)**(n-s)
like /= like.max() * 3
ax.plot(theta_grid, like, label="Likelihood (norm.)",
        linestyle=":")
ax.axvline(theta_true, color="red", alpha=0.5,
           label=f"True value ({theta_true})")
ax.set_xlabel(r"$\theta$")
ax.set_ylabel("Density")
ax.legend()
ax.set_title("Bayesian update: Bernoulli-Beta")
plt.tight_layout()
plt.savefig("bayes_update.pdf")
\end{minted}
\end{codeblock}

% ------------------------------------------------------------
\section{Exercises}
% ------------------------------------------------------------

\begin{exercise}
Let $X_1, \ldots, X_n \overset{\text{iid}}{\sim} \text{Poi}(\lambda)$
with prior $\lambda \sim \text{Ga}(\alpha, \beta)$.
\begin{enumerate}
  \item Determine the posterior distribution of $\lambda$.
  \item Compute the posterior mean and variance.
  \item Show that the posterior mean is a weighted average of the prior
        mean and the sample mean.
\end{enumerate}
\end{exercise}

\begin{exercise}
Compute the Jeffreys prior for the $\text{Ber}(\theta)$ model and show
that it is $\text{Be}(1/2, 1/2)$.
\end{exercise}

\begin{exercise}
Prove that in the Normal--Normal model, the posterior precision equals
the sum of the prior precision and the data precision. Interpret this
result.
\end{exercise}

\begin{exercise}
Consider two experiments:
\begin{enumerate}
  \item A coin is tossed $n = 12$ times, yielding $s = 3$ successes.
  \item A coin is tossed until the third success, requiring $n = 12$ tosses.
\end{enumerate}
Show that the likelihoods are proportional. Deduce that Bayesian
inference is identical, whereas the frequentist $p$-values differ.
\end{exercise}

\begin{exercise}
\textbf{(Numerical)} Consider the Normal--Normal model with
$\sigma^2 = 1$, $\mu_0 = 0$, $\tau_0^2 = 10$, and a sample of size
$n = 5$ drawn from $\mathcal{N}(2, 1)$.
\begin{enumerate}
  \item Compute the posterior analytically.
  \item Verify with PyMC.
  \item Plot the prior, posterior, and normalized likelihood on the
        same graph.
\end{enumerate}
\end{exercise}
